\section{Methodology for Model Optimization}\label{sec:method}

This section describes the self-recorded dataset, the data processing pipeline, the hyperparameter tuning of the four classifiers, and the performance evaluation method used to optimize the models on our 3-member group dataset.

\subsection{Dataset and Environment}
\begin{itemize}[leftmargin=*]
    \item \textbf{Data Source}: \texttt{bci\_dataset\_115-1}, containing self-recorded data from our 3 group members (30 recordings per subject, 90 in total).
    \item \textbf{Signal Type}: Single-channel (Fp1) EEG time-domain signal recorded via BrainLink with a dry electrode.
    \item \textbf{Sampling Rate}: 512 Hz; each recording covers one 20~s block (10240 samples).
    \item \textbf{Labels}: Each subject has ten rounds of recordings for three states (\texttt{\_1\_} = Relax, \texttt{\_2\_} = Focus, \texttt{\_3\_} = Stroop).
    \item \textbf{Environment}: Python 3.11 managed via \texttt{uv}, utilizing Scikit-learn, NumPy, SciPy, and other required packages.
\end{itemize}

\subsection{Data Preprocessing and Feature Engineering}
\todo{Describe in detail all the processing steps you took to improve model performance.}
\begin{itemize}[leftmargin=*]
    \item \textbf{Data Segmentation:}
    \begin{itemize}
        \item \texttt{SEGMENT\_LENGTH}: \todo{State the segment length} (in seconds) you chose and justify your choice, considering the trade-off between frequency resolution and the number of training samples.
        \item \texttt{OVERLAP\_RATIO}: \todo{State the overlap ratio} you chose and explain why.
    \end{itemize}

    \item \textbf{Preprocessing Pipeline:}
    \todo{List your processing steps in order.}
    \begin{enumerate}
        \item \todo{Step 1: }
    \end{enumerate}

    \item \textbf{Feature Engineering:}
    \todo{Explain how you transformed the raw signal into meaningful features for the classifiers.}
    \begin{itemize}
        \item \textbf{Frequency-Domain Features}: \todo{Explain how you extracted band power for the three states. Relate your choice to the expected physiology.}
        \item \textbf{Time-Domain Features}: \todo{Explain if you extracted statistical features and what they contribute beyond the spectral features.}
        \item \textbf{Feature Scaling}: \todo{State whether features were standardized.}
        \item \textbf{Feature Selection:}
        \begin{itemize}
            \item \texttt{FEATURE\_SELECTION}: \todo{State} whether you enabled feature selection (True/False).
            \item \texttt{N\_FEATURES\_SELECT}: \todo{If enabled, state} the number of features you selected and the basis for this choice.
        \end{itemize}
    \end{itemize}
\end{itemize}

\subsection{Classifiers and Hyperparameter Tuning}
This experiment requires implementing and comparing at least four classifiers: \texttt{MLPClassifier}, \texttt{SVC}, \texttt{RandomForestClassifier}, and one self-selected model. \todo{State which model you selected as the fourth classifier and why.} The table below lists the final settings used in this report; \todo{describe how these settings were obtained and report what you observed while tuning.}

\begin{table}[!htb]
    \centering
    \caption{Final Hyperparameter Settings of the Four Classifiers}
    \label{tab:hyperparams}
    \begin{tabularx}{\textwidth}{@{} l l l X @{}}
        \toprule
        \textbf{Model} & \textbf{Parameter} & \textbf{Final Setting} & \textbf{Justification for this Setting} \\
        \midrule
        \multirow{3}{*}{MLP} & \texttt{HIDDEN\_LAYER\_SIZES} & \todo{} & \todo{Briefly explain why.} \\
                             & \texttt{LEARNING\_RATE\_INIT} & \todo{} & \todo{Briefly explain why.} \\
                             & \texttt{ALPHA} (L2) & \todo{} & \todo{Briefly explain why.} \\
        \midrule
        \multirow{3}{*}{SVM} & \texttt{kernel} & \todo{} & \todo{Briefly explain why.} \\
                             & \texttt{C} & \todo{} & \todo{Briefly explain why.} \\
                             & \texttt{gamma} & \todo{} & \todo{Briefly explain why.} \\
        \midrule
        \multirow{2}{*}{RF}  & \texttt{n\_estimators} & \todo{} & \todo{Briefly explain why.} \\
                             & \texttt{max\_depth} & \todo{} & \todo{Briefly explain why.} \\
        \midrule
        \todo{Self-selected} & \todo{} & \todo{} & \todo{Briefly explain why.} \\
        \bottomrule
    \end{tabularx}
\end{table}

\subsection{Post-processing}
\todo{Describe whether you used any post-processing methods. If not, state "Not used."}
\begin{itemize}[leftmargin=*]
    \item \textbf{Probability Thresholds / Majority Voting:} \todo{If used, describe how utilizing \texttt{predict\_proba} or a sliding window helped smooth the predictions and reduce misclassifications of single segments, and state whether the decision was made per segment or aggregated over a whole block.}
\end{itemize}

\subsection{Performance Evaluation Method}
All four classifiers are evaluated with \textbf{LOSO Cross-Validation} to assess their generalization ability across different users. In this method, the data of one group member is used as the test set, while the data from the remaining 2 group members are used for training. This process is repeated 3 times, with each member serving as the test set once, and the reported result is the mean accuracy with its standard deviation across the three folds. In addition to overall accuracy, the per-class precision and recall and the $3\times3$ confusion matrix of the best model are reported in Section~\ref{sec:result}.